% ═══════════════════════════════════════════════════════════════════
% ЧАСТЬ X. РУКОВОДСТВО ПО ВОСПРОИЗВЕДЕНИЮ
% ═══════════════════════════════════════════════════════════════════
\part{Руководство по воспроизведению}

\section{Соответствие: раздел монографии $\leftrightarrow$ пункт
лаборатории}\label{sec:labmap}

Каждый численный результат монографии имеет прямой источник в меню
лаборатории \texttt{laboratory.py}. Таблица соответствия:

\begin{center}
\small
\begin{tabularx}{\textwidth}{@{}l l >{\raggedright\arraybackslash}X@{}}
\toprule
Раздел монографии & Числа & Пункт лаборатории \\
\midrule
\ref{sec:torus} (тор) & $(4,1,4\pi^2)$; $\gamma$; $\Delta$ &
«Стенды $\to$ Тор» \\
\ref{sec:k3} (K3) & $\rho=20$; $(1,19)$; $64$ & «Стенды $\to$ K3» \\
\ref{sec:e8} (errata) & $28/36$; $64$ & «Стенды $\to$ Errata E8» \\
\ref{sec:binary} (код) & $48/212/432/114$ & «Стенды $\to$ Бинарный
код» \\
\ref{sec:certA}--\ref{sec:certB} & ядро $29$; сборка $[Z]$ &
«Сертификаты $\to$ A, B» \\
\ref{sec:certC}--\ref{sec:certD} & $22=1+7+7+7$; $[L,P]=0$ &
«Сертификаты $\to$ C, D» \\
\ref{sec:certE} & $t^*=48$ & «Сертификаты $\to$ E» \\
\ref{sec:certF} & $Q=256$; $4X^2-7$ & «Сертификаты $\to$ F» \\
\ref{sec:certG} & Коула--Селберг; $\pi/15$, $\pi/30$ & «Сертификаты
$\to$ G» \\
\ref{sec:klein} (Клейн) & $j=-3375$; $\Omega$ & «Стенды $\to$
Клейн» \\
\ref{sec:cycles}--\ref{sec:indices} & V1--V9; SNF; $\mathrm{vol}_h$
& «Стенд N=15/30» (полный) \\
\ref{sec:certH} (H) & H1--H6 & «Сертификаты $\to$ H» \\
\bottomrule
\end{tabularx}
\end{center}

\subsection{Быстрый старт воспроизведения}

\begin{enumerate}[leftmargin=2em, itemsep=0.15em]
\item Установить зависимости: \texttt{pip install mpmath numpy
  matplotlib}.
\item Запустить лабораторию: \texttt{python3 laboratory.py}.
\item Выбрать язык интерфейса (русский/английский).
\item Пункт «Полный протокол V1--V9» --- воспроизводит таблицу
  \ref{tab:vprotocol} (~3 минуты).
\item Пункт «Отчёты и диаграммы» --- создаёт \texttt{reports/}:
  JSON-отчёты, текстовый журнал, плиточные диаграммы $600$~dpi.
\item Пункт «Мультиязычная верификация» --- при наличии
  \texttt{julia}, \texttt{gfortran/clang}, \texttt{cargo},
  \texttt{lean} запускает параллельные реализации и сверяет вердикты.
\end{enumerate}

\subsection{Правила сравнения результатов}

При сравнении чисел монографии с выводом лаборатории действует
правило порога: значение считается воспроизведённым, если
относительное отклонение не превышает заявленного порога протокола
($10^{-30}$ для $dps=35$, $10^{-65}$ для $dps=70$, целые числа ---
точное равенство). Пороги зашиты в отчёты; лаборатория печатает
вердикт PASS/FAIL по каждой проверке автоматически, и сводный вердикт
строки «все проверки PASS» эквивалентен воспроизведению раздела.

\section{Верификация на пяти языках и ядро Lean}\label{sec:multilang}

\subsection{Роли языков}

\begin{center}
\small
\begin{tabularx}{\textwidth}{@{}l l >{\raggedright\arraybackslash}X@{}}
\toprule
Язык & Файл & Роль \\
\midrule
Python & \texttt{laboratory.py} & полная лаборатория: протоколы,
отчёты, графики \\
Julia & \texttt{verification/julia/verify\_hodge.jl} &
высокоточные периоды и фазы \\
Fortran & \texttt{verification/fortran/verify\_hodge.f90} &
арифметика переписей и рангов \\
C & \texttt{verification/c/verify\_hodge.c} & быстрые целочисленные
проверки (Бареисс) \\
Rust & \texttt{verification/rust/verify\_hodge.rs} & безопасная
целочисленная арифметика \\
Lean~4 & \texttt{verification/lean/HodgeLaboratory.lean} & точные
тождества в ядре \\
\bottomrule
\end{tabularx}
\end{center}

Каждая реализация проверяет один и тот же базовый набор тождеств:
роды $91$ и $406$; переписи $1+6+84$ и $1+6+9+30+84+276$;
дискриминант $1125^2$; SNF-произведение; $b_{Ch}$ радикалы;
завершимость $t^*=\lcm$; сверки $48/212/432/114$; Клейн $j=-3375$.
Согласие пяти независимых реализаций и ядра Lean --- сильнейшая форма
сертификата воспроизводимости: ошибка компилятора или библиотеки,
повторённая в пяти экосистемах одновременно, практически исключена.

\subsection{Lean: что именно доказывает ядро}

Файл \texttt{HodgeLaboratory.lean} содержит формализованные
предложения, проверяемые ядром Lean~4 без аксиом сверх стандартных:
\begin{itemize}[leftmargin=1.6em, itemsep=0.1em]
\item род кривой Ферма: $(15-1)*(15-2)/2=91$,
  $(30-1)*(30-2)/2=406$;
\item переписи по проводникам: точные суммы Мёбиуса
  $1+6+84=91$, $1+6+9+30+84+276=406$;
\item дискриминант стенда $3^4*5^6=1265625=1125^2$;
\item произведение SNF-типа $(1,1,5,5,15,15,15,15)$ равно
  $1265625$;
\item радикальные значения $b_{Ch}(15)$ и $b_{Ch}(30)$ как
  выражения в $\sqrt5$, $\sqrt3$, $\sqrt{10\pm2\sqrt5}$;
\item завершимость: $t^*=\lcm(W/g, H/g)$ и частные случаи
  $t^*=48$ для сетки $48\times48$;
\item сверка бинарного кода: $48,212,432,114$;
\item целочисленные инварианты Клейна: $105^3=343*3375$,
  $j=-3375=-15^3$.
\end{itemize}
Ядро Lean проверяет эти предложения полным перебором вычислений
(\texttt{native\_decide}, \texttt{norm\_num}, \texttt{ring}) ---
это машинная проверка в самом строгом смысле; запуск выполняется
пунктом меню «Lean-верификация» либо командой
\texttt{lean HodgeLaboratory.lean}.

\subsection{Установка окружений (опционально)}

Все пять дополнительных окружений опциональны: лаборатория
функционирует полностью на Python; при наличии окружений она
автоматически расширяет вердикты. Краткие команды установки:

\begin{itemize}[leftmargin=1.6em, itemsep=0.1em]
\item Julia: \texttt{juliaup} или пакетный менеджер дистрибутива;
  запуск: \texttt{julia verification/julia/verify\_hodge.jl}.
\item Fortran: \texttt{gfortran -O2 -o verify\_hodge
  verify\_hodge.f90 \&\& ./verify\_hodge}.
\item C: \texttt{gcc -O2 -o verify\_hodge verify\_hodge.c -lm
  \&\& ./verify\_hodge}.
\item Rust: \texttt{rustc -O verification/rust/verify\_hodge.rs -o
  verify\_hodge \&\& ./verify\_hodge} (или \texttt{cargo run}).
\item Lean: \texttt{elan} (пакет \texttt{leanprover/lean4:v4.x});
  запуск: \texttt{lean verification/lean/HodgeLaboratory.lean}.
\end{itemize}

\section{Чек-лист полного воспроизведения}\label{sec:checklist}

\begin{enumerate}[leftmargin=2em, itemsep=0.15em]
\item \texttt{python3 laboratory.py} $\to$ выбрать «Полный
  протокол» $\to$ сводный вердикт PASS.
\item Отчёты в \texttt{reports/}: JSON-файлы совпадают с
  эталонными в \texttt{results/} по ключам и порогам.
\item Диаграммы \texttt{reports/plots/}: восемь плиток $600$~dpi
  сгенерированы без ошибок.
\item (Опционально) пять языков: все вердикты PASS, счётчики
  совпадают.
\item (Опционально) Lean: файл компилируется без ошибок,
  все \texttt{example} приняты ядром.
\item GitHub Actions (после пуша): workflow \texttt{ci.yml} зелёный
  --- прогон Python-протокола на трёх версиях Python.
\end{enumerate}

После прохождения чек-листа монография полностью воспроизведена:
каждое число таблиц \ref{tab:vprotocol}--\ref{tab:certs} и всех
разделов подтверждено независимым прогоном.
